\subsection{Problem 3. Trigonometric modeling}
The number of hours of daylight $D$ in a certain city can be approximated as a function of the day of the year $t$ (measured in days, with $t=0$ on January 1) by a sinusoidal model:
\[D(t)=A+B \cos(\omega(t-t_{0})).\]
Suppose you are given that the longest day has 14 hours of daylight (around day $t\approx172$) and the shortest day has 10 hours of daylight (around day $t\approx355$).

\begin{enumerate}[label=(\alph*)]
    \item Determine $A$ and $B$ from the information about the longest and shortest days. Explain the meaning of these two parameters.
    \item Argue that a reasonable choice for $\omega$ is $\dfrac{2\pi}{365}$. Explain briefly where this comes from.
    \item Choose a suitable value of $t_{0}$ so that the maximum of $D(t)$ occurs around $t=172$. Write the explicit function $D(t)$ with your chosen parameters.
    \item Using your model, estimate the number of hours of daylight on day $t=80$ and discuss qualitatively how accurate you expect the model to be (e.g., weather effects, geographical assumptions).
\end{enumerate}

\subsection*{Strategy}
\subsection*{Solution}

\pagebreak